For $x$ in a homogeneous basis of $\mathfrak{psl}(4|4)$, $\lambda_x$ is its $\mathrm{ad}_\Delta$ eigenvalue, $-1$ on $\Sigma$, $+1$ on $\Sigma^*$ and $+2$ for $K_b$, and the degree of $x$ in the position is $0$, $1$ and $2$ (Section 10.9). The two-body kernel of a pair of letters along $x$ is the kernel of Proposition \ref{prop:c10-arity1}; this section computes it along the supercharges, the special supercharges and the special conformal fields. Along $\Sigma$ the vertex $S^{(1)}(\xi)$ meets the gauge-fixing operator in a new way: the commutator $[\xi.,Q^{GF}]$ joins the gaugino to the ghost, and $S^{(1)}(\xi)$ sits on a cycle as a line from $\lambda$ to $c$ whose image under $Q$ is nonzero, in contrast to the composite lines along $\Delta$ and $P_b$. The section shows that the component along $x$ of every local functional is nevertheless a sum over its pairs of letters of local two-body kernels, the finite parts of the cycles with one vertex $S^{(1)}(x)$ and one special edge, plus a one-body term; writes these kernels as simplex integrals, independent of the endpoint along $K_b$; computes the action-level obstructions and solutions along the odd generators, the rotations and $K_b$ on two legs with all their free directions, and the arity-two equations that relate the free coefficients; assembles $Q_1$, $S_1$, $D_1$, $P_1$, $M_1$, $R_1$ and $K_1$ on the scalar bilinears, where the two-body part of $Q_1$ on the traceless part is $-2\log\tfrac32$ units times the classical transformation, as representatives with their scheme-independent content; and evaluates the pair entries of the twisted differential along a rank-one supercharge. Edges carry $P(\varepsilon,\infty)$ along $\Sigma$ and $\Sigma^*$ (Proposition \ref{prop:c8-Tsafe}) and $P(\varepsilon,T)$ at a finite $T$ along $K_b$ (Definition \ref{def:c8-K}). In this section $\psi$ denotes the finite part of the operators of Definition \ref{def:c8-renorm}, and $L:=\log\tfrac32$.
